\section{模型假设}

\begin{enumerate}[leftmargin=2.4em, itemsep=3pt]
  \item \textbf{连续介质假设}：药材视为连续介质，温度与含水率在空间上连续分布，
        可用偏微分方程描述。
  \item \textbf{轴对称假设}：药材为圆柱，几何与边界条件轴对称，
        故采用\textbf{一维轴对称径向模型}；端面效应以二维轴对称对照验证。
  \item \textbf{局部热平衡假设}：同一空间点上固相与液相温度相同，
        不需要分别写两相能量方程。
  \item \textbf{均匀仿射收缩假设（仅问题4）}：固定长度 $+$ 均匀径向仿射收缩，
        材料径向速度 $v_r=(\dot R/R)\,r$。
        \textbf{必须声明这是假设而非结论}——附件2 只给整体半径 $R(t)$，
        无内部位移数据，单个表面半径不能识别内部位移是否仿射。
  \item \textbf{物性经验公式适用性}：$\rho$、$c_p$、$k$、$D$ 一律采用附录给出的
        经验公式，\textbf{不替换文献参数}，也不引入题面未给的物性。
  \item \textbf{环境边界等效约定}：按题设等效约定，把附件1 的``水分浓度''列
        作为药材表面的\textbf{等效平衡含水率} $\Cinf$，
        \textbf{不假装}它是由真实空气状态求出的吸附平衡结果。
  \item \textbf{一维径向近似的合理性}：端面方向与径向的扩散时间尺度比约 $39$，
        端面影响由二维轴对称对照定量验证。
  \item \textbf{模型分层}：M0 为\textbf{题设基线（四问正式答案）}；
        M1 为可核查扩展（潜热、辐射、端面、结壳、参数情景），条件充分时做；
        M2 为工程推广，只给路线与数据需求。
\end{enumerate}

\subsection{入口检查（六项，逐条落实）}

为避免``用了题面之外的量''这类不可追溯的错误，建模前先做六项入口检查：

\begin{table}[H]
  \centering
  \caption{入口检查结论表}
  \label{tab:entry}
  \footnotesize
  \begin{tabularx}{\textwidth}{@{}clX@{}}
    \toprule
    \# & 检查项 & 结论 \\
    \midrule
    1 & 含水率 $C$ 与空气含湿量 $Y$ 的基准 &
      $C=m_w/m_d$（干基）、$Y=m_v/m_{da}$；\textbf{即使同为 kg/kg 也不可自动相减}，
      本文引入已明确基准的 $\Cinf$ 表示等效环境值 \\
    2 & 密度 $\rho$ 的定义 &
      附录3/4 给的是\textbf{湿物料体积密度}，用于热容项；
      $\rho_d=\rho/(1+C)$ 与之不兼容（见 \S\ref{sec:geo-density}） \\
    3 & 温度单位 &
      内部一律摄氏度，\textbf{仅在 $D$ 公式内}转开尔文，转换点全局唯一 \\
    4 & 长时边界外推方式 &
      附件1 只到 14400 s，恒温段\textbf{必须显式假设}，不得用插值函数外推 \\
    5 & 输出口径 &
      严格按附件3 模板：表头含反斜杠、距离轴为数值、
      \textbf{result4 末列为字符串``药材表面''且距离轴只到 1.9 cm} \\
    6 & 达标判据 &
      题面``各处''$\Rightarrow$ 全域最大值 $\Cmax(t)$，
      \textbf{不得用平均值或表面值} \\
    \bottomrule
  \end{tabularx}
\end{table}
